\section[Demonstrations]{\secstyle{Demonstrations:}}

Soit \(u, v\) deux vecteurs d'un \(\R\)-espace vectoriel \(F\).

(1) Montrer que \(Vect(3u, 5v) = Vect(u, v)\) en utilisant uniquement le lemme d'expansion.

\begin{proof}[\unskip\nopunct] 
   \fbox{\(\subseteq\)} 
   On remarque que: 

   \begin{align*}
      3u = 3u + 0v \in Vect(u, v)\\ 
      5v = 0u + 5v \in Vect(u, v) 
   \end{align*}   
   On voit donc que \(\{3u, 5v\} \subseteq Vect(u, v)\) et par le lemme d'expansion, on conclut que \(Vect(3u, 5v) \subseteq Vect(u, v)\).\<

   \fbox{\(\supseteq\)}
   On remarque que: 
   \begin{align*}
      u &= \frac{1}{3}(3u) + 0v \in Vect(3u, 5v)\\ 
      v &= 0u + \frac{1}{5}5v \in Vect(3u, 5v) 
   \end{align*}   
   On voit donc que \(\{u, v\} \subseteq Vect(3u, 5v)\) et par le lemme d'expansion, on conclut que \(Vect(u, v) \subseteq Vect(3u, 5v)\).

\end{proof}


\section[Exercices]{\secstyle{Exercices:}}


